\documentclass[10pt,a4paper]{amsart}
\usepackage[margin=19mm]{geometry}
\usepackage{amsmath,amssymb,mathtools}
\usepackage[scr=rsfs,cal=euler]{mathalfa}
\usepackage{xfrac}
\usepackage[unicode,hidelinks,bookmarks=false]{hyperref}
\newcommand{\ie}{i.e.\ }
\newcommand{\cf}{cf.\ }
\newcommand{\resp}{respectively\ }
\newcommand{\Subsection}{\subsection}
\providecommand{\nobreakdash}{\nobreak\hbox{-}}
\setlength{\emergencystretch}{2em}
\multlinegap0pt
\usepackage{bm}
\usepackage{bbm}
\usepackage{dsfont}
\usepackage{xspace}
\usepackage{cancel}
\usepackage{mathtools}
\usepackage{amsthm}
\makeatletter
\@ifundefined{thm}{\newtheorem{thm}{\bf Theorem}}{}
\@ifundefined{lem}{\newtheorem{lem}[thm]{\bf Lemma}}{}
\makeatother
\DeclareMathOperator{\diam}{diam}
\DeclareMathOperator{\dist}{dist}
\DeclareMathOperator{\supp}{supp}
\newcommand{\mm}{\mathfrak m}
\title[Associated primes of the second power of closed neighborhood]{Associated primes of the second power of closed neighborhood ideals of graphs}
\author{Ha Thi Thu Hien, Thanh Vu}
\date{Source: arXiv:2507.08777v1 (CC BY 4.0)}
\begin{document}
\maketitle

\noindent\emph{Source and licence.} Associated primes of the second power of closed neighborhood ideals of graphs. Version arXiv:2507.08777v1 (CC BY 4.0), distributed under the Creative Commons Attribution 4.0 International licence, as recorded in the arXiv licence metadata for this exact version and shown on the abstract page https://arxiv.org/abs/2507.08777v1. Full rights notice: licenses/arxiv-2507.08777-CC-BY-4.0.txt.\par\smallskip

\noindent\textit{Editorial scope.} Counted unit: the Lem whose source label is \texttt{lem\_directed\_cycle} in the article (source0.tex lines 419--426, source label \texttt{lem\_directed\_cycle}), delivered together with the article's own complete original proof of it (source0.tex lines 428--442, one \texttt{proof} environment of 15 lines). No mathematics is written, repaired or re-derived: statement and proof are the article's own text line for line. Required context retained uncounted: thm\_diam\_2 (lines 299--301); thm\_diam\_3 (lines 306--316). Every definition, display and label of those lines is retained unchanged. Omitted from this package and not redistributed: 1--298, 302--305, 317--418, 427--427, 443--666. Nothing inside a retained excerpt is omitted or altered: every retained body is delivered exactly as the article's lines are written, so no omission receipt of any kind accompanies this package. Citations: the retained excerpts carry no citation command at all, so no bibliography entry of the article is transcribed and no reference is redistributed. Reference adaptation: none was needed and none was made; every reference inside the delivered unit resolves inside it, so no unresolved reference or citation is produced. Printed slips kept literal: no printed word, symbol, hypothesis or number of the counted statement or of its proof was altered. Layout and class: the article's own class is \texttt{amsart} with packages including amssymb, amsmath, amscd, enumerate, verbatim, hyperref, inputenc, color, graphicx, tikz, xy, pstricks, pst-3d; the wrapper is the lane's pinned \texttt{amsart} editorial wrapper at 10pt on A4, which restates the theorem-like environments the retained text needs and adds the article's own macro definitions that the retained text needs (\texttt{\textbackslash diam}, \texttt{\textbackslash dist}, \texttt{\textbackslash mm}, \texttt{\textbackslash supp}), with extra wrapper packages (bm, bbm, dsfont, xspace, cancel, mathtools). The delivered numbering is the unit's own. Assets: no retained span contains a figure environment, a graphics-inclusion command, a PDF-page inclusion, a table or a TikZ picture; no image file is copied into this package, the package declares no asset dependency and no image file is redistributed with it. Licence: version arXiv:2507.08777v1 of this article is distributed under the Creative Commons Attribution 4.0 International licence, as recorded in the arXiv licence metadata for this exact version and shown on the abstract page https://arxiv.org/abs/2507.08777v1. A full rights notice is delivered at licenses/arxiv-2507.08777-CC-BY-4.0.txt. Characters: every retained excerpt is pure ASCII, so the delivered engine, XeLaTeX with its default Latin Modern text font, reports no missing character.
\begin{thm}\label{thm_diam_2}
Let $G$ be a simple graph of diameter at most $2$, and let $NI(G)$ denote the closed neighborhood ideal of $G$. Then $\mm$ is an associated prime of $S/NI(G)^2$ if and only if $G$ is vertex diameter-$2$-critical.
\end{thm}

\begin{thm}\label{thm_diam_3}
Let $G$ be a simple graph with $\diam(G) \geq 3$, and let $NI(G)$ be its closed neighborhood ideal. Let $H$ be the hypergraph associated to $G$ as described above. Then $\mm$ is an associated prime of $S/NI(G)^2$ if and only if there exists a subset $C \subseteq V(G)$ such that:
\begin{enumerate}
    \item $C$ is a vertex cover of $H$;
    \item for every $i \in C$, there exists an edge $e \in E(H)$ such that $C \cap e = \{i\}$;
    \item for every $j \in V(G) \setminus C$, there exist vertices $j_1, j_2 \in N_G(j)$ such that $N_G[j_1], N_G[j_2] \subseteq V(G) \setminus C$, and 
    \[
    \dist_G(j_1, j_2) = 2, \quad \text{and} \quad \dist_{G - j}(j_1, j_2) \geq 3.
    \]
\end{enumerate}
\end{thm}

\begin{lem}\label{lem_directed_cycle}
With the notation above, we have:
\begin{enumerate}
    \item $W \neq \emptyset$;
    \item $\dist_G(i,j) \le 2$ for all $i,j \in W$;
    \item the induced subgraph of $G$ on $W$ contains a cycle.
\end{enumerate}
\end{lem}

\begin{proof}
(i) Since $U = \supp (f)$ where $\mm = NI(G)^2 : f$ for some squarefree monomial $f$, we deduce that $U$ is non-empty. By Theorems~\ref{thm_diam_2} and~\ref{thm_diam_3}, for any $u \in U$, there exist $v_1, v_2 \in N_G(u)$ such that $N_G[v_1]$ and $N_G[v_2]$ are subsets of $U$. By the definition of $W$, this means $v_1, v_2 \in W$. Thus, $W \neq \emptyset$.

\vspace{1em} % Adds a bit of vertical space between items

(ii) Assume for the sake of contradiction that $\dist_G(i,j) \ge 3$ for some $i,j \in W$. This implies that their closed neighborhoods are disjoint, i.e., $N_G[i] \cap N_G[j] = \emptyset$. Since $i,j \in W$, their closed neighborhoods $N_G[i]$ and $N_G[j]$ are both subsets of $U$. Therefore, $N_G[i] \cup N_G[j] \subseteq U$. This leads to a contradiction, as $C = V(G) \setminus U$ does not cover the edge $N_G[i] \cup N_G[j]$ of $H$.

\vspace{1em}

(iii) To establish this, we construct a directed graph $D$ with vertex set $W$. By Theorems~\ref{thm_diam_2} and~\ref{thm_diam_3}, for each $u \in W \subseteq U$, there exist $v_1, v_2 \in W$ such that $\dist_G(v_1, v_2) = 2$ and $\dist_{G - u}(v_1, v_2) \geq 3$. For each $u$, we fix one such pair $(v_1, v_2)$ and add directed edges $(u, v_1)$ and $(u, v_2)$ to $D$. Note that $v_1, v_2 \in N_G(u)$, so the underlying undirected graph of $D$ is a subgraph of the induced subgraph of $G$ on $W$.

By construction, every vertex in $D$ has out-degree exactly 2. Since $W$ is finite, following directed edges from any starting vertex eventually leads to a repeated vertex, which implies the existence of a directed cycle in $D$.

If any such directed cycle has length at least $3$, we are done. Suppose instead that all directed cycles in $D$ have length $2$. Then we may remove such a cycle and repeat the argument on the remaining vertices. If at any stage we find a longer cycle, we are done. Otherwise, we conclude that $D$ is a disjoint union of directed $2$-cycles. In this case, the underlying undirected graph of $D$ is a graph in which every vertex has degree exactly 2, implying that it is a disjoint union of (undirected) cycles. Therefore, the induced subgraph of $G$ on $W$ contains a cycle, completing the proof.
\end{proof}

\end{document}
